\section{Passive selected-port cut}
\label{sec:passive}

Consider first a finite-dimensional energy-normalized passive realization
\begin{equation}
A=-iH-\frac12K^\dagger K,
\qquad H=H^\dagger,
\label{eq:Apassive}
\end{equation}
where $K$ contains every physical radiation and loss port. Let $K_i$ and $K_o$ select two disjoint port groups with no direct cross-feedthrough. Their strictly proper transfer block is
\begin{equation}
H_{o\leftarrow i}(\nu)
=-K_o(i\nu I-A)^{-1}K_i^\dagger .
\label{eq:crossH}
\end{equation}
A selected local drive therefore excites endpoint dynamics before gravitational radiation is emitted, and a separated receiver is driven by the retarded TT field. Models containing a direct selected-input-to-selected-output bypass lie outside the present $H_2$ statement.

Define the selected-input Gramian
\begin{equation}
P_i(\tau)=\int_0^\tau
e^{At}K_i^\dagger K_i e^{A^\dagger t}\dd t .
\end{equation}
Since $K_i^\dagger K_i\le K^\dagger K=-(A+A^\dagger)$,
\begin{equation}
\boxed{
0\le P_i(\tau)
\le I-e^{A\tau}e^{A^\dagger\tau}
\le I .
}
\label{eq:PiBound}
\end{equation}
For a stable realization, Plancherel gives
\begin{align}
\|H_{o\leftarrow i}\|_2^2
&\equiv \frac{1}{2\pi}\int_{-\infty}^{\infty}\dd\nu\,
\Tr(H_{o\leftarrow i}^\dagger H_{o\leftarrow i})\\
&=\Tr(K_oP_iK_o^\dagger)
\le \Tr(K_o^\dagger K_o).
\end{align}
The dual output Gramian gives the complementary input bound, so
\begin{equation}
\boxed{
\|H_{o\leftarrow i}\|_2^2
\le
\min\!\left[
\Tr(K_i^\dagger K_i),
\Tr(K_o^\dagger K_o)
\right].
}
\label{eq:h2cut}
\end{equation}
This is generic passive-system mathematics, not a gravitational result.

A weighted form of the same inequality will be useful below. For any bounded map $L$ acting on the endpoint state,
\begin{equation}
\frac{1}{2\pi}\int_{-\infty}^{\infty}\dd\nu\,
\left\|L(i\nu I-A)^{-1}K_i^\dagger\right\|_{\HS}^2
=\Tr(LP_iL^\dagger)
\le\Tr(L^\dagger L).
\label{eq:weightedcut}
\end{equation}
No assumption that $L$ itself is a physical loss port is needed; Eq.~\eqref{eq:weightedcut} follows directly from $P_i\le I$.

For the separated gravitational link retain the physical propagation frequency rather than freezing it at the carrier:
\begin{equation}
T(\nu)=H_B(\nu)P_g[\omega(\nu),R]H_A(\nu).
\label{eq:factor}
\end{equation}
Here $H_A$ maps selected source inputs to the retained gravitational output space, $H_B$ maps retained gravitational input channels to selected receiver outputs, and $P_g$ is the free-space propagation operator. Pointwise passivity gives $\|H_A(\nu)\|_{\op},\|H_B(\nu)\|_{\op}\le1$.

For later comparison define
\begin{equation}
\bar\eta(R)
=\sup_{\nu\in\mathcal B_\nu}
\|P_g[\omega(\nu),R]\|_{\op}^2 .
\end{equation}
Then positivity of the integrand, extension of the band integral to the full line only for the endpoint $H_2$ block, and Eq.~\eqref{eq:h2cut} give
\begin{equation}
\Gamma_{\rm coh}
\le
\bar\eta(R)\Tr(K_{g,A}^\dagger K_{g,A}),
\end{equation}
and independently
\begin{equation}
\Gamma_{\rm coh}
\le
\bar\eta(R)\Tr(K_{g,B}^\dagger K_{g,B}).
\end{equation}
Thus the familiar scalar cut is
\begin{equation}
\boxed{
\Gamma_{\rm coh}
\le
\bar\eta(R)
\min\!\left[
\Tr(K_{g,A}^\dagger K_{g,A}),
\Tr(K_{g,B}^\dagger K_{g,B})
\right].
}
\label{eq:gravcut}
\end{equation}
Unlike a carrier-frozen factorization, Eq.~\eqref{eq:gravcut} already retains the propagation variation across the measured band.

More importantly, Eq.~\eqref{eq:weightedcut} permits the propagation operator to weight different gravitational angular sectors before the endpoint trace is taken. Section~\ref{sec:propagation} uses this stronger matrix form to avoid multiplying a total quadrupole-strength ceiling by an independently optimized directivity ceiling.

The same Gramian inequalities extend to a separable countably infinite modal Hilbert space when the Markov port operators are bounded and the retained gravitational port is Hilbert--Schmidt. The retained resource in Sec.~\ref{sec:resource} supplies the required finite trace. Arbitrary unbounded boundary-control PDE ports require a separate admissibility analysis \cite{BarasBrockett1975,Opmeer2013}.
